\documentclass[10pt,a4paper]{amsart}
\usepackage{amsmath,amssymb,amsthm,graphicx}
% For commutative diagrams we use Xy-pic
\usepackage[all]{xy}

% We use 2cell for 2-commutative diagrams.
\xyoption{2cell}
% Theorem environments.
%
\theoremstyle{plain}
\newtheorem{theorem}[subsection]{Theorem}
\newtheorem{proposition}[subsection]{Proposition}
\newtheorem{lemma}[subsection]{Lemma}

\theoremstyle{definition}
\newtheorem{definition}[subsection]{Definition}
\newtheorem{example}[subsection]{Example}
\newtheorem{exercise}[subsection]{Exercise}
\newtheorem{situation}[subsection]{Situation}

\theoremstyle{remark}
\newtheorem{remark}[subsection]{Remark}
\newtheorem{remarks}[subsection]{Remarks}

\numberwithin{equation}{subsection}

% Macros
%
\def\lim{\mathop{\mathrm{lim}}\nolimits}
\def\colim{\mathop{\mathrm{colim}}\nolimits}
\def\Spec{\mathop{\mathrm{Spec}}}
\def\Hom{\mathop{\mathrm{Hom}}\nolimits}
\def\Ext{\mathop{\mathrm{Ext}}\nolimits}
\def\SheafHom{\mathop{\mathcal{H}\!\mathit{om}}\nolimits}
\def\SheafExt{\mathop{\mathcal{E}\!\mathit{xt}}\nolimits}
\def\Sch{\mathit{Sch}}
\def\Mor{\mathop{\mathrm{Mor}}\nolimits}
\def\Ob{\mathop{\mathrm{Ob}}\nolimits}
\def\Sh{\mathop{\mathit{Sh}}\nolimits}
\def\NL{\mathop{N\!L}\nolimits}
\def\CH{\mathop{\mathrm{CH}}\nolimits}
\def\proetale{{pro\text{-}\acute{e}tale}}
\def\etale{{\acute{e}tale}}
\def\QCoh{\mathit{QCoh}}
\def\Ker{\mathop{\mathrm{Ker}}}
\def\Im{\mathop{\mathrm{Im}}}
\def\Coker{\mathop{\mathrm{Coker}}}
\def\Coim{\mathop{\mathrm{Coim}}}

% Boxtimes
%
\DeclareMathSymbol{\boxtimes}{\mathbin}{AMSa}{"02}

%
% Macros for moduli stacks/spaces
%
\def\QCohstack{\mathcal{QC}\!\mathit{oh}}
\def\Cohstack{\mathcal{C}\!\mathit{oh}}
\def\Spacesstack{\mathcal{S}\!\mathit{paces}}
\def\Quotfunctor{\mathrm{Quot}}
\def\Hilbfunctor{\mathrm{Hilb}}
\def\Curvesstack{\mathcal{C}\!\mathit{urves}}
\def\Polarizedstack{\mathcal{P}\!\mathit{olarized}}
\def\Complexesstack{\mathcal{C}\!\mathit{omplexes}}
% \Pic is the operator that assigns to X its picard group, usage \Pic(X)
% \Picardstack_{X/B} denotes the Picard stack of X over B
% \Picardfunctor_{X/B} denotes the Picard functor of X over B
\def\Pic{\mathop{\mathrm{Pic}}\nolimits}
\def\Picardstack{\mathcal{P}\!\mathit{ic}}
\def\Picardfunctor{\mathrm{Pic}}
\def\Deformationcategory{\mathcal{D}\!\mathit{ef}}
\usepackage{cleveref}
\crefname{originallemma}{original lemma}{original lemmas}
\Crefname{originallemma}{Original lemma}{Original lemmas}
\makeatletter
\newlabel{algebra-lemma-normal-goes-up}{{033C}{source}{}{sourcepointer-1}{}}
\newlabel{algebra-lemma-normal-goes-up@cref}{{[originallemma][0][]033C}{[1][0][]source}{}{}{}}
\makeatother
\crefname{originallemma}{original lemma}{original lemmas}
\Crefname{originallemma}{Original lemma}{Original lemmas}
\makeatletter
\newlabel{algebra-lemma-normal-ring-integrally-closed}{{034M}{source}{}{sourcepointer-2}{}}
\newlabel{algebra-lemma-normal-ring-integrally-closed@cref}{{[originallemma][0][]034M}{[1][0][]source}{}{}{}}
\makeatother
\crefname{originaltheorem}{original theorem}{original theorems}
\Crefname{originaltheorem}{Original theorem}{Original theorems}
\makeatletter
\newlabel{etale-theorem-etale-radicial-open}{{025G}{source}{}{sourcepointer-3}{}}
\newlabel{etale-theorem-etale-radicial-open@cref}{{[originaltheorem][0][]025G}{[1][0][]source}{}{}{}}
\makeatother
\crefname{originalreference}{original reference}{original references}
\Crefname{originalreference}{Original reference}{Original references}
\makeatletter
\newlabel{lemma-describe-etale-local-ring}{{04HX}{source}{}{sourcepointer-4}{}}
\newlabel{lemma-describe-etale-local-ring@cref}{{[originalreference][0][]04HX}{[1][0][]source}{}{}{}}
\makeatother
\crefname{originallemma}{original lemma}{original lemmas}
\Crefname{originallemma}{Original lemma}{Original lemmas}
\makeatletter
\newlabel{morphisms-lemma-distinct-local-rings}{{02NF}{source}{}{sourcepointer-5}{}}
\newlabel{morphisms-lemma-distinct-local-rings@cref}{{[originallemma][0][]02NF}{[1][0][]source}{}{}{}}
\makeatother
\crefname{originallemma}{original lemma}{original lemmas}
\Crefname{originallemma}{Original lemma}{Original lemmas}
\makeatletter
\newlabel{morphisms-lemma-universally-injective}{{01S4}{source}{}{sourcepointer-6}{}}
\newlabel{morphisms-lemma-universally-injective@cref}{{[originallemma][0][]01S4}{[1][0][]source}{}{}{}}
\makeatother
\crefname{originallemma}{original lemma}{original lemmas}
\Crefname{originallemma}{Original lemma}{Original lemmas}
\makeatletter
\newlabel{pione-lemma-normal-local-domain-separablly-closed-fraction-field}{{0BSQ}{source}{}{sourcepointer-7}{}}
\newlabel{pione-lemma-normal-local-domain-separablly-closed-fraction-field@cref}{{[originallemma][0][]0BSQ}{[1][0][]source}{}{}{}}
\makeatother
\title{On an integral normal scheme with separably closed function field every separated etale morphism splits as a disjoint union of open immersions and all local rings are strictly henselian}
\author{The Stacks Project Authors}
\date{stacks-project:etale-cohomology.tex:local-snapshot-20261011; GNU FDL 1.2 or later}
\begin{document}
\maketitle
\noindent\textit{Editorial note.} All delivered mathematical text is the author's own, copied from the pinned Stacks Project chapter file etale-cohomology.tex; no statement or proof step was rewritten, reordered or removed.\par\medskip
\noindent\textit{Editorial note.} The counted proof is the author's complete argument: the proof is written out in full: the affine splitting argument with the idempotents and the radicial criterion, the global decomposition along the specialisations of the generic points of the fibre with the separatedness step, and the deduction of part (2) from the description of strict henselisations.\par\medskip
\noindent\textit{Editorial note.} The label lemma-normal-scheme-with-alg-closed-function-field names the counted statement, and the proof environment immediately following it in the source is delivered verbatim as its proof.\par\medskip
\noindent\textit{Editorial note.} The counted statement is carried by the author's own lemma environment in the chapter source.\par\medskip
\noindent\textit{Editorial note.} The statement and proof are copied from the pinned chapter file; the seven cited inputs are kept as external pointers with their canonical Stacks tags.\par\medskip
\noindent\textit{Editorial note.} The preamble is the chapter's own native preamble excerpt already used for this source: the Xy-pic setup (native preamble lines 14-18) and the native theorem declarations with the mathematical macros (native preamble lines 154-217).\par\medskip
\noindent\textit{Allowed standard prerequisites (editorial).}\begin{itemize}
\item Etale ring maps and their behaviour after tensoring with the fraction field. Same chapter, the structure of an etale map R -> A along a normal domain R with separably closed fraction field. Supplies the splitting A tensor\_R K = product of copies of K and the idempotents e\_i used at the start of the proof.
\item Normality passes up along etale maps and integrally closed domains have no new idempotents. Algebra, the lemmas that normality goes up and that a normal domain is integrally closed inside its fraction field. Supplies the fact that the idempotents e\_i lie in A, giving the product decomposition A = product A e\_i.
\item Criterion for an etale radicial map to be an open immersion. Etale Morphisms, the theorem characterising an etale universally injective morphism as an open immersion, via the criterion of Morphisms. Supplies the conclusion that R -> A defines an open immersion once A tensor\_R A -> A is bijective.
\item Strict henselian local rings as local rings of the etale site. Same chapter, the description of the strict henselisation of O\_\{X,x\} as the local ring at a geometric point on the etale site. Supplies part (2) from part (1).
\end{itemize}\medskip
\ifdefined\hypertarget\hypertarget{sourcepointer-1}{}\fi
\noindent\textit{Source reference.} Original lemma 033C: sources/stacks-chow/tags.txt, source lines 4051--4051. Prerequisite pointer invoked at proof line(s) 14328; it names a standard fact of a sibling chapter that the counted argument uses, and every step of the counted proof built on it is retained verbatim in the delivered body.\par\medskip
\ifdefined\hypertarget\hypertarget{sourcepointer-2}{}\fi
\noindent\textit{Source reference.} Original lemma 034M: sources/stacks-chow/tags.txt, source lines 4097--4097. Prerequisite pointer invoked at proof line(s) 14330; it names a standard fact of a sibling chapter that the counted argument uses, and every step of the counted proof built on it is retained verbatim in the delivered body.\par\medskip
\ifdefined\hypertarget\hypertarget{sourcepointer-3}{}\fi
\noindent\textit{Source reference.} Original theorem 025G: sources/stacks-chow/tags.txt, source lines 2813--2813. Prerequisite pointer invoked at proof line(s) 14341; it names a standard fact of a sibling chapter that the counted argument uses, and every step of the counted proof built on it is retained verbatim in the delivered body.\par\medskip
\ifdefined\hypertarget\hypertarget{sourcepointer-4}{}\fi
\noindent\textit{Source reference.} Original reference 04HX: sources/stacks-etale-cohomology/pinned-original-source.tex, source lines 4257--4257. Prerequisite pointer invoked at proof line(s) 14362; it names a standard fact of the same chapter that the counted argument uses, and every step of the counted proof built on it is retained verbatim in the delivered body.\par\medskip
\ifdefined\hypertarget\hypertarget{sourcepointer-5}{}\fi
\noindent\textit{Source reference.} Original lemma 02NF: sources/stacks-chow/tags.txt, source lines 3478--3478. Prerequisite pointer invoked at proof line(s) 14354; it names a standard fact of a sibling chapter that the counted argument uses, and every step of the counted proof built on it is retained verbatim in the delivered body.\par\medskip
\ifdefined\hypertarget\hypertarget{sourcepointer-6}{}\fi
\noindent\textit{Source reference.} Original lemma 01S4: sources/stacks-chow/tags.txt, source lines 2325--2325. Prerequisite pointer invoked at proof line(s) 14339; it names a standard fact of a sibling chapter that the counted argument uses, and every step of the counted proof built on it is retained verbatim in the delivered body.\par\medskip
\ifdefined\hypertarget\hypertarget{sourcepointer-7}{}\fi
\noindent\textit{Source reference.} Original lemma 0BSQ: sources/stacks-chow/tags.txt, source lines 14813--14813. Prerequisite pointer invoked at proof line(s) 14365; it names a standard fact of a sibling chapter that the counted argument uses, and every step of the counted proof built on it is retained verbatim in the delivered body.\par\medskip
\begin{lemma}
\label{lemma-normal-scheme-with-alg-closed-function-field}
Let $X$ be an integral normal scheme with separably closed
function field.
\begin{enumerate}
\item A separated \'etale morphism $U \to X$ is a
disjoint union of open immersions.
\item All local rings of $X$ are strictly henselian.
\end{enumerate}
\end{lemma}

\begin{proof}
Let $R$ be a normal domain whose fraction field is separably algebraically
closed. Let $R \to A$ be an \'etale ring map. Then
$A \otimes_R K$ is as a $K$-algebra a finite product
$\prod_{i = 1, \ldots, n} K$ of copies of $K$. Let $e_i$, $i = 1, \ldots, n$
be the corresponding idempotents of $A \otimes_R K$. Since $A$ is normal
(Algebra, Lemma \ref{algebra-lemma-normal-goes-up})
the idempotents $e_i$ are in $A$
(Algebra, Lemma \ref{algebra-lemma-normal-ring-integrally-closed}).
Hence $A = \prod Ae_i$ and we may assume $A \otimes_R K = K$.
Since $A \subset A \otimes_R K = K$ (by flatness of $R \to A$ and
since $R \subset K$) we conclude that $A$ is a domain.
By the same argument we conclude that
$A \otimes_R A \subset (A \otimes_R A) \otimes_R K = K$.
It follows that the map $A \otimes_R A \to A$ is
injective as well as surjective. Thus $R \to A$ defines an
open immersion by
Morphisms, Lemma \ref{morphisms-lemma-universally-injective}
and
\'Etale Morphisms, Theorem \ref{etale-theorem-etale-radicial-open}.

\medskip\noindent
Let $f : U \to X$ be a separated \'etale morphism. Let $\eta \in X$
be the generic point and let $f^{-1}(\{\eta\}) = \{\xi_i\}_{i \in I}$.
The result of the previous paragraph shows the following:
For any affine open $U' \subset U$ whose image in $X$ is contained in
an affine we have $U' = \coprod_{i \in I} U'_i$ where $U'_i$
is the set of point of $U'$ which are specializations of $\xi_i$.
Moreover, the morphism $U'_i \to X$ is an open immersion.
It follows that $U_i = \overline{\{\xi_i\}}$ is an open and closed
subscheme of $U$ and that $U_i \to X$ is locally on the source
an isomorphism. By Morphisms,
Lemma \ref{morphisms-lemma-distinct-local-rings}
the fact that $U_i \to X$ is separated, implies that
$U_i \to X$ is injective and we conclude that $U_i \to X$
is an open immersion, i.e., (1) holds.

\medskip\noindent
Part (2) follows from part (1) and the description of the strict
henselization of $\mathcal{O}_{X, x}$ as the local ring at $\overline{x}$
on the \'etale site of $X$ (Lemma \ref{lemma-describe-etale-local-ring}).
It can also be proved directly, see
Fundamental Groups, Lemma
\ref{pione-lemma-normal-local-domain-separablly-closed-fraction-field}.
\end{proof}

\end{document}
